\par\medskip
\noindent
1.2. Apliquemos VIII~7.1(b) al par $(A,A')$ y a la biextensión
$W_K$ de $(A_K,A'_K)$ por $\mathbb G_{mK}$ de (1.0.1). Obtenemos así
un \underline{emparejamiento canónico}
\begin{equation}\tag{1.2.1}
 \Phi_o\times\Phi'_o\longrightarrow(\mathbb Q/\mathbb Z)_k,
\end{equation}
es decir, un emparejamiento de $\pi_o$-módulos
\begin{equation}\tag{1.2.2}
 \Phi_o(\overline k)\times\Phi'_o(\overline k)
 \longrightarrow\mathbb Q/\mathbb Z.
\end{equation}
Aquí el emparejamiento aparece como la obstrucción a prolongar $W_K$
a una biextensión $W$ de $(A,A')$ por $\mathbb G_{mS}$; el significado
de este punto de vista se explica en VIII~7.3.2. Enunciamos de inmediato
la siguiente conjetura acerca de este emparejamiento.

\medskip
\noindent
\underline{Conjetura 1.3.} \underline{El emparejamiento} (1.2.1)
\underline{es una dualidad perfecta}: \underline{identifica cada uno de
los grupos de componentes} $\Phi_o$ \underline{y} $\Phi'_o$
\underline{con el dual del otro, con valores en}
$(\mathbb Q/\mathbb Z)_k$.

\medskip
\noindent
1.3.1. La formación de (1.2.1) conmuta con el cambio de base
formalmente étale (VIII~7.3.5.3), y la formación de los modelos de
Néron posee la propiedad análoga. Por consiguiente, para demostrar la
Conjetura~1.3 bastaría tratar el caso en que $S$ es completo y posee
cuerpo residual separablemente cerrado. Artin y Mazur demostraron la
conjetura cuando $A_K$ es la jacobiana de una curva propia lisa sobre
$K$, mediante su teoría general, inédita, de la autodualidad de la
jacobiana relativa (en el sentido de las categorías derivadas) de un
esquema de curvas propio sobre $S$. También demostraron la Conjetura~1.3
cuando el cuerpo residual es \underline{finito}, mediante la dualidad
de Tate [32]. Veremos asimismo en 11.4 y 11.5 (a reserva de comprobaciones
de compatibilidad) que 1.3 es válida sobre las componentes
$\ell$-primarias
para $\ell\ne p$, y que lo es en el caso de «reducción semiestable».
Estos resultados hacen extremadamente plausible 1.3. Por otra parte,
1.3 puede considerarse también un subproducto de una teoría general de
dualidad para esquemas en grupos sobre cuerpos de fracciones de anillos
de valoración discreta completos, formalmente análoga a SGA~5~I. Tal
teoría sigue siendo por ahora heurística; contendría también la teoría
de cuerpos de clases local en la forma que le dio Serre [27].

\medskip
\noindent
1.3.2. Para un esquema abeliano dado $A_K$, la validez de 1.3 significa
que el emparejamiento (1.2.1) es \underline{separante} tanto por la
izquierda como por la derecha. Según VIII~7.3.2, la separación por la
izquierda expresa precisamente que el único subesquema en grupos
abierto $V$ de $A$ para el cual $W_K$ puede prolongarse a una
biextensión de $(V,A')$ por $\mathbb G_{mS}$ es la componente neutra
$A^o$ de $A$. Existe una interpretación simétrica de la separación
por la derecha.

\medskip
\noindent
1.4. Más generalmente, supongamos que
\begin{equation}\tag{1.4.1}
 U\subset A\qquad\text{y}\qquad U'\subset A'
\end{equation}
son subesquemas en grupos abiertos tales que $W_K$ se prolonga a una
biextensión de $(U,U')$ por $\mathbb G_{mS}$. Según VIII~7.3.2, esto
significa explícitamente que los subgrupos
\begin{equation}\tag{1.4.2}
 \sisi{\Phi'_o}{F_o}=U_o/U_o^o\subset\Phi_o,
 \qquad
 \sisi{\Psi'_o}{F'_o}=U'_o/{U'}_o^o
 \subset\sisi{\Psi_o}{\Phi'_o}
\end{equation}
definidos por $U$ y $U'$ son ortogonales respecto de (1.2.1).\footnote{La
fórmula impresa, confirmada mediante una comprobación directa de alta
resolución, reutiliza $\Phi'_o$ para el primer subgrupo y escribe
$\Psi'_o\subset\Psi_o$ para el segundo, entrando en conflicto con la
notación ya establecida para el grupo de componentes de $A'$. Los
nombres locales neutros $F_o,F'_o$ hacen inequívocos los dos subgrupos
pretendidos sin alterar el enunciado.} Este prolongamiento (único salvo
isomorfismo único) de $W_K$ se denota por $W$. Queda
determinado salvo isomorfismo único por su clase
\begin{equation}\tag{1.4.3}
 w\in\mathrm{Biext}^1(U,U';\mathbb G_{mS}),
\end{equation}
caracterizada a su vez por la condición de que su imagen en
$\mathrm{Biext}^1(A_K,A'_K;\mathbb G_{mK})$ sea la clase $w_K$ de
(1.0.1).

Por ejemplo, pueden tomarse como $U$ y $U'$ las componentes neutras
$A^o$ y $A'{}^o$ de $A$ y $A'$. Se obtiene así una
\underline{biextensión canónica $W^o$, o $W$, de $(A^o,A'{}^o)$ por
$\mathbb G_{mS}$}, que da lugar a una clase
\begin{equation}\tag{1.4.4}
 w=w^o\in\mathrm{Biext}^1(A^o,A'{}^o;\mathbb G_{mS}),
\end{equation}
con la que trabajaremos principalmente. Es evidente que las
restricciones a $(A^o,A'{}^o)$ de las clases $w$ de (1.4.3), para los
distintos pares $(U,U')$, son todas iguales a $w^o$ de (1.4.4).

\medskip
\noindent
\underline{Proposición 1.5.} \underline{Considérese la fibra especial
$W_s$ de la biextensión canónica} $W$ \underline{de}
$(A^o,A'{}^o)$ \underline{por} $\mathbb G_{mS}$ \underline{de (1.4.4),
de modo que $W_s$ es una biextensión de $(A_o^o,{A'}_o^o)$ por}
$\mathbb G_{mk}$. \underline{Esta biextensión es separante}
(VIII~4.11).

En efecto, elijamos un haz invertible amplio $\mathbb L$ sobre $A^o$.
Entonces $\mathbb L_\eta$ define un homomorfismo (cf. 1.0.5)
\[
 A_K\longrightarrow A'_K,
\]
que se prolonga a un homomorfismo $A\longrightarrow A'$ y, por tanto,
da
\[
 \psi:A^o\longrightarrow A'{}^o.
\]
Podemos considerar la biextensión $W_{\mathbb L}$ de $(A^o,A^o)$ por
$\mathbb G_{mS}$ obtenida al traer $W$ por
$(\operatorname{id}_{A^o},\psi)$. Su restricción a las fibras genéricas
es precisamente la biextensión $\delta(\mathbb L_\eta)$ de VIII~4.12.
Por tanto, según VIII~7.4, $W_{\mathbb L}$ está asociada al haz
invertible birigidificado $\delta(\mathbb L)$ y, en consecuencia,
$(W_{\mathbb L})_s$ está definida por $\delta(\mathbb L_s)$. Por otra
parte, puede describirse también como la imagen inversa de $W_s$ por
los morfismos
\[
 \operatorname{id}_{A_o^o}:A_o^o\longrightarrow A_o^o,
 \qquad
 \psi_o:A_o^o\longrightarrow {A'}_o^o.
\]
Puesto que $\mathbb L$, y por tanto $\mathbb L_s$, es amplio,
VIII~4.13 muestra que $W_{\mathbb L_s}$ es separante, en particular
separante por la izquierda. Con mayor razón, $W_s$ es por ello
separante por la izquierda. Por simetría, lo es también por la derecha,
lo que demuestra 1.5.

Reforzaremos sustancialmente 1.5 cuando $A_o^o$ tenga rango unipotente
cero (5.4).
